% Figure: Michaelis-Menten saturation and its Lineweaver-Burk linearisation.
% (a) rate versus substrate, saturating; (b) double-reciprocal straight line.
\begin{figure}[htbp]
\centering
\begin{tikzpicture}[font=\small]
% ---- panel (a): saturation curve ----
\begin{scope}
\draw[->, thick, draw=engblack!70] (0,0) -- (5.2,0)
  node[right, font=\scriptsize] {$[S]$};
\draw[->, thick, draw=engblack!70] (0,0) -- (0,4.0)
  node[above, font=\scriptsize] {$v$};
% v = Vmax S/(Km+S), Vmax=3.4 (drawn), Km at half-max
\draw[very thick, draw=engblue!80] plot[domain=0:5,samples=60]
  (\x, {3.4*\x/(1.2+\x)});
% Vmax asymptote
\draw[dashed, draw=enggray] (0,3.4) -- (5,3.4);
\node[font=\scriptsize, color=enggray, anchor=east] at (5.0,3.55) {$V_{\max}$};
% half-max line and Km
\draw[dashed, draw=engred!70] (0,1.7) -- (1.2,1.7) -- (1.2,0);
\node[font=\scriptsize, color=engred!70, anchor=north] at (1.2,-0.05) {$K_m$};
\node[font=\scriptsize, anchor=north] at (2.6,-0.5) {(a) saturation};
% data points
\foreach \px in {0.3,0.7,1.2,1.9,2.8,4.0} {
  \pgfmathsetmacro\vy{3.4*\px/(1.2+\px) + 0.09*sin(\px*300)}
  \filldraw[engblack!75] (\px, {\vy}) circle (0.05);
}
\end{scope}
% ---- panel (b): Lineweaver-Burk ----
\begin{scope}[xshift=7.2cm]
\draw[->, thick, draw=engblack!70] (-1.4,0) -- (5.2,0)
  node[right, font=\scriptsize] {$1/[S]$};
\draw[->, thick, draw=engblack!70] (0,-0.2) -- (0,4.0)
  node[above, font=\scriptsize] {$1/v$};
% 1/v = (Km/Vmax)(1/S) + 1/Vmax; slope Km/Vmax, intercept 1/Vmax.
% choose intercept 0.6, slope 0.55; x-intercept = -1/Km
\draw[very thick, draw=engblue!80] (-1.1, 0.0) -- (5.0, 3.36);
% y-intercept marker
\filldraw[engred!80] (0,0.605) circle (0.05);
\node[font=\scriptsize, color=engred!80, anchor=west] at (0.15,0.75) {$1/V_{\max}$};
% x-intercept marker
\filldraw[engred!80] (-1.1,0) circle (0.05);
\node[font=\scriptsize, color=engred!80, anchor=south] at (-1.1,0.1) {$-1/K_m$};
% slope label
\node[font=\scriptsize, color=engblue!80, anchor=north west, rotate=28] at (2.5,1.7)
  {slope $=K_m/V_{\max}$};
% data points on the line
\foreach \px in {-0.9,-0.4,0.4,1.3,2.4,3.7} {
  \pgfmathsetmacro\vy{0.605+0.55*\px + 0.06*sin(\px*400)}
  \filldraw[engblack!75] (\px, {\vy}) circle (0.05);
}
\node[font=\scriptsize, anchor=north] at (2.0,-0.5) {(b) double reciprocal};
\end{scope}
\end{tikzpicture}
\caption[Michaelis-Menten kinetics and the Lineweaver-Burk
linearisation]{Enzyme reaction rate against substrate concentration. Panel
(a): the rate $v = V_{\max}[S]/(K_m+[S])$ rises toward the saturation
$V_{\max}$, with the half-saturation constant $K_m$ read where the rate
reaches half its ceiling. Panel (b): the double-reciprocal transform
$1/v = (K_m/V_{\max})(1/[S]) + 1/V_{\max}$ turns the saturating curve into
a straight line whose intercepts identify both parameters. The
transformation is a worked instance of the chapter's linearise-to-identify
move, with the working caution that the reciprocal reweights the noise:
low-substrate points, where $1/[S]$ is large, dominate the fit, so the
graphical read gives starting values while the honest interval comes from
a nonlinear fit on the untransformed rate.}
\label{fig:vol02ch18:michaelis-menten}
\end{figure}
